\section{Algebraic intersections and KKT as reference routes}
\label{sec:algebraic}

\subsection{Why tangency appears as a repeated root}

In the projective chart, write the loss circle and voltage conic as
\begin{equation}
 G_c(\eta,\chi)=\eta^2+\chi^2-c\eta+1=0,
 \qquad F(\eta,\chi)=0.
\end{equation}
For a fixed \(c\), eliminating \(\chi\) with
\(p(\eta;c)=\operatorname{res}_\chi(F,G_c)\) gives a polynomial whose roots
are the projected intersection coordinates. Two quadratic conics produce a
quartic intersection polynomial in the generic case. Before first contact
there are no real roots; at contact two intersection points coalesce; after
contact they separate. Coalescence is precisely a multiple root, so
\begin{equation}
 p(\eta;c)=0,\qquad \frac{\partial p}{\partial\eta}(\eta;c)=0,
 \quad\Longleftrightarrow\quad \disc_\eta p(\eta;c)=0.
 \label{eq:discriminant-contact}
\end{equation}
This is the algebraic version of common tangent lines.

For the EESM's generic projective conic, fixed-level intersection is quartic.
If the level parameter is eliminated as well, direct stationarity of
\(J(\eta,\chi)\) on \(F=0\) combines one quadratic and one cubic equation;
the univariate stationary equation is generically degree six. Axis symmetry,
zero resistance, or vanishing cross terms can factor it into lower-degree
pieces, frequently quartics or polynomials in a squared variable. Thus
``quadrics'' does not imply that every EESM contact is solved by one universal
quartic. Candidate roots must still satisfy \(\eta>0\), the original voltage
inequality, the requested torque sign, and the unsquared equations.

For the 2D PSM, intersection of the current circle, voltage ellipse, torque
hyperbola, or MTPV locus likewise produces quartics. Closed radical formulas
exist and have been developed systematically for anisotropic synchronous
machines including resistance and cross-coupling \cite{eldeeb2018,eldeeb2016}.
They are valuable offline and for exact regression. On a controller, a
monotone bracketed scalar search often has better coefficient conditioning,
clearer branch selection, and predictable execution than a general quartic in
radicals.

\subsection{Lagrange and KKT after the geometry is visible}

When voltage is inactive, first contact between a loss ellipsoid and torque
surface has parallel normals. This gives
\begin{equation}
 \gls{sym:rmat}\gls{sym:i}=\lambda\gls{sym:hmat}\gls{sym:i},
 \label{eq:loss-torque-gep}
\end{equation}
the generalized eigenproblem already explained by loss whitening. Its largest
positive generalized eigenvalue gives the best torque per loss. The equation
is not sufficient by itself: the desired torque fixes scale, and the voltage
inequality must be checked.

With voltage potentially active, minimum loss at fixed torque satisfies
\begin{align}
 (\matr R+\lambda\matr H+\mu\matr Q_U)\vect i&=\vect0,
 \label{eq:minloss-kkt}\\
 \vect i^{\mathsf T}\matr H\vect i&=M_{\mathrm{ref}},\\
 \vect i^{\mathsf T}\matr Q_U\vect i&\leq U_{\max}^2,\qquad
 \mu\geq0,\qquad
 \mu(\vect i^{\mathsf T}\matr Q_U\vect i-U_{\max}^2)=0.
\end{align}
The multiplier \(\mu\) is zero if the voltage cylinder is not touched. If it
is active, the loss normal is balanced by torque and voltage normals. The
nonconvex torque equality means a KKT point is only a candidate; globality
comes here from comparing branches or, preferably, from the convex
performance-space boundary.

Maximum torque under both inequalities has
\begin{equation}
 (-\matr H+\mu_P\matr R+\mu_U\matr Q_U)\vect i=\vect0,
 \quad \mu_P,\mu_U\geq0,
 \label{eq:mmax-kkt}
\end{equation}
with two complementary-slackness equations. This active-set language remains
useful for verifying whether the loss ellipsoid, voltage cylinder, or both
limit a returned point. It is also the natural extension when one adds field
or stator-current bounds as in real EESM setpoint algorithms
\cite{reinhard2022}.

\paragraph{Controller judgment.}
Resultants provide finite candidate sets and symbolic insight; KKT provides
residuals and active-set labels. Neither is selected as the primary online
coordinate because the performance boundary reduces branch management before
forming high-degree polynomials.
